% GENERATED FILE. Edit canonical lessons, then run npm run generate:books.
% canonical-source-hash: fnv1a64:d0bf2ebda3490a12
% canonical-lessons: RU-C314-salyut, RU-C314-feyerverk, RU-C314-konfetti, RU-C314-khlopushka, RU-C314-sharik

\chapter{Fireworks}
\label{ch:ru-fireworks}
\hlchaptermodality{314}

\begin{chapteropening}
\textbf{By the end of this chapter:} \emph{I can talk about a firework display, fireworks, confetti, a party popper and a balloon.}

Complete the last lesson of chapter 314: I can talk about a firework display, fireworks, confetti, a party popper and a balloon.
\end{chapteropening}

\section[\texorpdfstring{\ru{салют} (salyút)}{салют (salyút)}]{\ru{салют} (salyút) — a firework display; a salute}

\label{lesson:RU-C314-salyut}



\begin{quote}
Before the new one: say the Russian for a decoration, then the Russian for to dress up.
\end{quote}

\subsection*{You'll want to know: \ru{салют}}
\textbf{\ru{салют}} (\emph{salyút}) — \textquotedblleft{}a firework display; a salute\textquotedblright{}.

\begin{quote}
\ru{В} \ru{полночь} \ru{начался} \ru{салют}.
\end{quote}

\begin{quote}
\emph{V pólnoch nachalsyá salyút.}
\end{quote}

\begin{quote}
At midnight the firework display began.
\end{quote}

\emph{\ru{Праздничный} \ru{салют}} — a festive firework display.

\begin{cousinweb}
\emph{\ru{салют}} is usually traced to French \emph{salut} \textquotedblleft{}greeting\textquotedblright{}, from Latin \emph{salūs} \textquotedblleft{}health, safety\textquotedblright{} (Lewis \& Short); EtymDB glosses the Russian word \textquotedblleft{}salute; fireworks\textquotedblright{} but lists no parent.
\end{cousinweb}

\begin{grammarlens}[title={What you've built}]
One more word for telling, explaining and describing at B1.
\end{grammarlens}

\subsection*{Your turn}
\begin{itemize}
\raggedright
  \item \emph{Say it:} \emph{salyút}
  \item \emph{Say it:} \emph{salyút}, once more
  \item \emph{Say it:} a sentence of your own with \emph{salyút}
  \item \emph{Recall:} say the Russian for a decoration, then the Russian for to dress up, then say \emph{salyút} again
\end{itemize}

\begin{tcolorbox}[breakable,colback=teal!4,colframe=teal!35!black,title={Before you move on}]
What does \ru{салют} mean? (\textquotedblleft{}A firework display; a salute\textquotedblright{}.) Say it once more.
\end{tcolorbox}
\section[\texorpdfstring{\ru{фейерверк} (feyervérk)}{фейерверк (feyervérk)}]{\ru{фейерверк} (feyervérk) — fireworks}

\label{lesson:RU-C314-feyerverk}



\begin{quote}
Before the new one: say the Russian for to dress up, then the Russian for a firework display.
\end{quote}

\subsection*{You'll want to know: \ru{фейерверк}}
\textbf{\ru{фейерверк}} (\emph{feyervérk}) — \textquotedblleft{}fireworks\textquotedblright{}.

\begin{quote}
\ru{Мы} \ru{смотрели} \ru{фейерверк} \ru{с} \ru{балкона}.
\end{quote}

\begin{quote}
\emph{My smotréli feyervérk s balkóna.}
\end{quote}

\begin{quote}
We watched the fireworks from the balcony.
\end{quote}

Masculine.

\begin{cousinweb}
\emph{\ru{фейерверк}} was borrowed from German \emph{Feuerwerk} (EtymDB), from \emph{Feuer} \textquotedblleft{}fire\textquotedblright{} and \emph{Werk} \textquotedblleft{}work\textquotedblright{}.
\end{cousinweb}

\begin{grammarlens}[title={What you've built}]
One more word for telling, explaining and describing at B1.
\end{grammarlens}

\subsection*{Your turn}
\begin{itemize}
\raggedright
  \item \emph{Say it:} \emph{feyervérk}
  \item \emph{Say it:} \emph{feyervérk}, once more
  \item \emph{Say it:} a sentence of your own with \emph{feyervérk}
  \item \emph{Recall:} say the Russian for to dress up, then the Russian for a firework display, then say \emph{feyervérk} again
\end{itemize}

\begin{tcolorbox}[breakable,colback=teal!4,colframe=teal!35!black,title={Before you move on}]
What does \ru{фейерверк} mean? (\textquotedblleft{}Fireworks\textquotedblright{}.) Say it once more.
\end{tcolorbox}
\section[\texorpdfstring{\ru{конфетти} (konfettí)}{конфетти (konfettí)}]{\ru{конфетти} (konfettí) — confetti}

\label{lesson:RU-C314-konfetti}



\begin{quote}
Before the new one: say the Russian for a firework display, then the Russian for fireworks.
\end{quote}

\subsection*{You'll want to know: \ru{конфетти}}
\textbf{\ru{конфетти}} (\emph{konfettí}) — \textquotedblleft{}confetti\textquotedblright{}.

\begin{quote}
\ru{Конфетти} \ru{было} \ru{везде}.
\end{quote}

\begin{quote}
\emph{Konfettí býlo vezdé.}
\end{quote}

\begin{quote}
There was confetti everywhere.
\end{quote}

Neuter and indeclinable.

\begin{cousinweb}
\emph{\ru{конфетти}} was borrowed from Italian \emph{confetti} (EtymDB), the plural of \emph{confetto} \textquotedblleft{}a sugared sweet\textquotedblright{}, once thrown at carnivals; it shares its source, Latin \emph{conficere} \textquotedblleft{}to prepare\textquotedblright{}, with \emph{\ru{конфета}}.
\end{cousinweb}

\begin{grammarlens}[title={What you've built}]
One more word for telling, explaining and describing at B1.
\end{grammarlens}

\subsection*{Your turn}
\begin{itemize}
\raggedright
  \item \emph{Say it:} \emph{konfettí}
  \item \emph{Say it:} \emph{konfettí}, once more
  \item \emph{Say it:} a sentence of your own with \emph{konfettí}
  \item \emph{Recall:} say the Russian for a firework display, then the Russian for fireworks, then say \emph{konfettí} again
\end{itemize}

\begin{tcolorbox}[breakable,colback=teal!4,colframe=teal!35!black,title={Before you move on}]
What does \ru{конфетти} mean? (\textquotedblleft{}Confetti\textquotedblright{}.) Say it once more.
\end{tcolorbox}
\section[\texorpdfstring{\ru{хлопушка} (khlopúshka)}{хлопушка (khlopúshka)}]{\ru{хлопушка} (khlopúshka) — a party popper, a cracker}

\label{lesson:RU-C314-khlopushka}



\begin{quote}
Before the new one: say the Russian for fireworks, then the Russian for confetti.
\end{quote}

\subsection*{You'll want to know: \ru{хлопушка}}
\textbf{\ru{хлопушка}} (\emph{khlopúshka}) — \textquotedblleft{}a party popper, a cracker\textquotedblright{}.

\begin{quote}
\ru{Хлопушка} \ru{громко} \ru{хлопнула}.
\end{quote}

\begin{quote}
\emph{Khlopúshka grómko khlópnula.}
\end{quote}

\begin{quote}
The party popper went off with a loud bang.
\end{quote}

Feminine; genitive plural \emph{\ru{хлопушек}}.

\begin{cousinweb}
\emph{\ru{хлопушка}} is built on \emph{\ru{хлопать}} \textquotedblleft{}to bang, to clap\textquotedblright{}, a word that imitates the sound; EtymDB lists no parent.
\end{cousinweb}

\begin{grammarlens}[title={What you've built}]
One more word for telling, explaining and describing at B1.
\end{grammarlens}

\subsection*{Your turn}
\begin{itemize}
\raggedright
  \item \emph{Say it:} \emph{khlopúshka}
  \item \emph{Say it:} \emph{khlopúshka}, once more
  \item \emph{Say it:} a sentence of your own with \emph{khlopúshka}
  \item \emph{Recall:} say the Russian for fireworks, then the Russian for confetti, then say \emph{khlopúshka} again
\end{itemize}

\begin{tcolorbox}[breakable,colback=teal!4,colframe=teal!35!black,title={Before you move on}]
What does \ru{хлопушка} mean? (\textquotedblleft{}A party popper, a cracker\textquotedblright{}.) Say it once more.
\end{tcolorbox}
\section[\texorpdfstring{\ru{шарик} (shárik)}{шарик (shárik)}]{\ru{шарик} (shárik) — a balloon; a little ball}

\label{lesson:RU-C314-sharik}



\begin{quote}
Before the new one: say the Russian for confetti, then the Russian for a party popper.
\end{quote}

\subsection*{You'll want to know: \ru{шарик}}
\textbf{\ru{шарик}} (\emph{shárik}) — \textquotedblleft{}a balloon; a little ball\textquotedblright{}.

\begin{quote}
\ru{Мальчик} \ru{держал} \ru{красный} \ru{шарик}.
\end{quote}

\begin{quote}
\emph{Málchik derzhál krásnyy shárik.}
\end{quote}

\begin{quote}
The boy was holding a red balloon.
\end{quote}

A diminutive of \emph{\ru{шар}} \textquotedblleft{}a ball, a sphere\textquotedblright{}; \emph{\ru{воздушный} \ru{шарик}} — a balloon.

\begin{cousinweb}
\emph{\ru{шарик}} is a diminutive of \emph{\ru{шар}} \textquotedblleft{}a ball\textquotedblright{}, whose origin is not settled; EtymDB lists no parent.
\end{cousinweb}

\begin{grammarlens}[title={What you've built}]
One more word for telling, explaining and describing at B1.
\end{grammarlens}

\subsection*{Your turn}
\begin{itemize}
\raggedright
  \item \emph{Say it:} \emph{shárik}
  \item \emph{Say it:} \emph{shárik}, once more
  \item \emph{Say it:} a sentence of your own with \emph{shárik}
  \item \emph{Recall:} say the Russian for confetti, then the Russian for a party popper, then say \emph{shárik} again
\end{itemize}

\begin{tcolorbox}[breakable,colback=teal!4,colframe=teal!35!black,title={Before you move on}]
What does \ru{шарик} mean? (\textquotedblleft{}A balloon; a little ball\textquotedblright{}.) Say it once more.
\end{tcolorbox}
